\par\noindent\textbf{(1)} 正则单粒子配分函数 z₁=V/λ\_T³，对不可分辨粒子计数 Z\_N=z₁\textasciicircum{}N/N!。用 Stirling lnN!≈NlnN−N 得 F=−k\_BT lnZ\_N≈−Nk\_BT[ln(V/(Nλ\_T³))+1]，此自由能为后续热力学导数的唯一基础，N!不可省略。
\par\noindent\textbf{(2)} 由 S=−(∂F/∂T)\_\{V,N\}，注意 λ\_T∝T\textasciicircum{}\{−1/2\}，得到 S=Nk\_B[ln(V/(Nλ\_T³))+5/2]；压力 P=−F\_V=Nk\_BT/V；化学势 μ=F\_N=k\_BT ln(nλ\_T³)。熵中5/2包括动量空间温度导数的3/2与广延性计数。
\par\noindent\textbf{(3)} 绝热固定S、N时 ln(VT\textasciicircum{}\{3/2\}/N)=常数，故 TV\textasciicircum{}\{2/3\}=常数，配合PV=Nk\_BT得 PV\textasciicircum{}\{5/3\}=常数，即单原子理想气体γ=5/3。此过程应是准静态可逆绝热，若自由膨胀不可逆产生熵则不能套同一关系。
\par\noindent\textbf{(4)} 若把Z\_N错误写成z₁\textasciicircum{}N，相同气体混合时会出现依赖于人为标签的虚假广延熵，且化学势失去正确密度依赖；N!消除该Gibbs悖论。高密低温时nλ\_T³不再小，Bose/Fermi统计替代Boltzmann配分函数，Stirling型Sackur–Tetrode不再适用。
